\documentclass[11pt,aspectratio=169]{beamer}

% Theme
\usetheme{Madrid}
\usecolortheme{default}

% Vietnamese support — use lualatex or xelatex for best results
% (works with pdflatex too via utf8 + T5, but T5 fontenc may need vntex)
\usepackage[utf8]{inputenc}
\usepackage[T5,T1]{fontenc}
\usepackage[main=vietnamese,english]{babel}

% Packages
\usepackage{booktabs}
\usepackage{graphicx}
\usepackage{tikz}
\usetikzlibrary{shadows,arrows.meta,positioning,calc}
\usepackage{amsmath}
\usepackage{xcolor}

% Royal Society palette
\definecolor{rsnavy}{HTML}{1B2A4A}
\definecolor{rsclaret}{HTML}{8E1A2B}
\definecolor{rsteal}{HTML}{16A085}
\definecolor{rsamber}{HTML}{F39C12}

% Larger default font for readability
\setbeamerfont{frametitle}{size=\Large,series=\bfseries}
\setbeamerfont{block title}{size=\normalsize,series=\bfseries}
\setbeamercolor{frametitle}{fg=rsnavy}
\setbeamercolor{title}{fg=rsnavy}
\setbeamercolor{section in toc}{fg=rsnavy}

% Custom take-away macro
\newcommand{\takeaway}[1]{%
  \begin{center}
    \colorbox{rsnavy!10}{%
      \begin{minipage}{0.92\textwidth}
        \centering\textbf{#1}
      \end{minipage}}
  \end{center}}

\title[Babel's Machines]{Babel's Machines}
\subtitle{Ngôn ngữ định hình lại hành vi\\của các tác tử nhân tạo đa ngôn ngữ}
\author[Đào S. D. Minh]{Đào Sỹ Duy Minh \and Vũ Nguyễn}
\institute[VNU-HCM]{Khoa Công nghệ Thông tin\\Trường Đại học Khoa học Tự nhiên, ĐHQG-HCM\\[4pt]
\textit{Cho Interface Focus theme issue về machine behaviour}}
\date{\today}

\begin{document}

% =====================================================
% Slide 1 — Tiêu đề
% =====================================================
\begin{frame}
    \titlepage
\end{frame}

% =====================================================
% Slide 2 — Roadmap
% =====================================================
\begin{frame}{Lộ trình}
    \large
    \begin{enumerate}
        \setlength\itemsep{0.5em}
        \item \textbf{Tại sao ngôn ngữ quan trọng} cho chương trình machine-behaviour
        \item \textbf{LLM agents làm gì} qua các ngôn ngữ — bao gồm cả mâu thuẫn
        \item \textbf{Tại sao} — cơ chế, alternatives, và vòng lặp causal cần đóng
        \item \textbf{Vậy thì sao} — tiến hóa, governance, và bốn dự đoán có thể bác bỏ
    \end{enumerate}
\end{frame}

% =====================================================
% PHẦN 1 — TẠI SAO NGÔN NGỮ QUAN TRỌNG
% =====================================================
\section{Tại sao ngôn ngữ quan trọng}

\begin{frame}{LLM không còn chỉ tạo văn bản. Chúng \emph{hành động}.}
    \large
    \begin{itemize}
        \setlength\itemsep{0.4em}
        \item Đàm phán giao dịch và tạo liên minh (Diplomacy, bargaining)
        \item Hợp tác và phản bội trong các trò chơi lặp lại (Prisoner's Dilemma)
        \item Lừa dối trong các trò chơi suy luận xã hội (Werewolf, Avalon)
        \item Quản lý ngân sách trong sàn đấu giá; cộng tác trên các tác vụ kéo dài hàng giờ
    \end{itemize}
    \vspace{0.4cm}
    \takeaway{Đây là \emph{ethology}, không phải benchmark. Công cụ đúng đến từ sinh học, lý thuyết trò chơi tiến hóa, và tiến hóa văn hóa --- không phải kỹ thuật.}
    \vspace{0.2cm}
    \begin{flushright}
        \small\itshape Rahwan et al., \emph{Nature} 2019 --- chương trình machine-behaviour
    \end{flushright}
\end{frame}

\begin{frame}{Luận điểm trong một câu}
    \vspace{0.2cm}
    \large
    \begin{quote}
        \itshape ``Ngôn ngữ mà LLM agent hoạt động trong đó định hình lại một cách hệ thống hành vi chiến lược, theory-of-mind, và prosocial defaults của nó.''
    \end{quote}
    \vspace{0.3cm}
    \normalsize
    Chúng tôi gọi đây là \textbf{Algorithmic Sapir--Whorf Switch} (theo Wang et al., 2025).
    \vspace{0.4cm}
    \begin{block}{Cái chúng tôi \emph{không} claim}
        Không phải claim Whorfian mạnh — rằng ngôn ngữ quyết định tư duy.\\
        Chỉ là: các biểu diễn đặc thù theo ngôn ngữ bên trong mô hình điều biến chính sách mà mô hình thực thi trong môi trường tương tác.\\
        \emph{Thực nghiệm, không siêu hình.} Chúng tôi dành không gian thực sự cho bằng chứng phản bác.
    \end{block}
\end{frame}

% =====================================================
% PHẦN 2 — LLM AGENTS LÀM GÌ
% =====================================================
\section{LLM agents làm gì qua các ngôn ngữ}

\begin{frame}{Bảy họ trò chơi tương tác --- tất cả đều cho thấy hiệu ứng ngôn ngữ}
    \large
    \begin{columns}[T]
        \begin{column}{0.48\textwidth}
            \begin{itemize}
                \setlength\itemsep{0.3em}
                \item Tình thế xã hội (IPD, Public Goods)
                \item Đàm phán (Diplomacy, bargaining)
                \item Thuyết phục / đấu giá
                \item Suy luận xã hội (Werewolf, Avalon, Spyfall)
            \end{itemize}
        \end{column}
        \begin{column}{0.48\textwidth}
            \begin{itemize}
                \setlength\itemsep{0.3em}
                \item Giao tiếp hợp tác (Codenames, Hanabi)
                \item Trò chơi mở thế giới (Minecraft, BALROG)
                \item Dialogue games (clembench)
            \end{itemize}
        \end{column}
    \end{columns}
    \vspace{0.5cm}
    \takeaway{Mọi họ đã được chạy đa ngôn ngữ. Mọi họ đều báo cáo divergence theo ngôn ngữ.}
\end{frame}

\begin{frame}{Cùng mô hình, cùng game, ngôn ngữ khác $\rightarrow$ chiến lược khác}
    \begin{columns}[T]
        \begin{column}{0.42\textwidth}
            \begin{table}
            \centering\small
            \begin{tabular}{@{}lcccc@{}}
                \toprule
                & ALLC & ALLD & TFT & WSLS \\
                \midrule
                Anh        & 26 & 25 & 25 & 24 \\
                Pháp       & 27 & 18 & 26 & 29 \\
                Ả Rập      & \textbf{32} & \textbf{28} & 19 & 21 \\
                Trung      & 31 & 30 & 19 & \textbf{20} \\
                Việt       & 26 & 24 & 23 & 27 \\
                \bottomrule
            \end{tabular}
            \end{table}
            \scriptsize Tổng hợp qua 4 họ mô hình,\\$\geq 200$ vòng/điều kiện.\\Tỷ lệ mô tả, không phải meta-analysis.
        \end{column}
        \begin{column}{0.55\textwidth}
            \textbf{Ba pattern cần chú ý:}
            \vspace{0.2cm}
            \begin{itemize}
                \setlength\itemsep{0.4em}
                \item \textcolor{rsclaret}{\textbf{Tiếng Ả Rập bimodal}} --- cả hợp tác cao nhất \emph{và} phản bội cao thứ hai. Volatile, không adversarial.
                \item \textbf{Tiếng Việt subverts collectivist prior} --- defection cao trong một số setup; một mô hình collapse xuống 2\% cooperation.
                \item \textbf{Tiếng Trung ưa Win-Stay-Lose-Shift} hơn Tit-for-Tat --- outcome-driven, không phải reciprocity-driven.
            \end{itemize}
            \vspace{0.2cm}
            \scriptsize\itshape Hiệu ứng họ mô hình cũng lấn át hiệu ứng ngôn ngữ: GPT-4o ``hawkish'', Mistral ``dovish'' --- bất kể ngôn ngữ.
        \end{column}
    \end{columns}
\end{frame}

\begin{frame}{LLM ``nicer than humans'' --- trong một regime rất cụ thể}
    \large
    LLM tiền tuyến hợp tác ở tỷ lệ cao hơn nhiều so với người chơi điển hình trong IPD lặp lại.
    \vspace{0.3cm}
    \begin{block}{Regime mà điều này đúng}
        Ma trận payoff hợp tác $(R{=}3, S{=}0, T{=}5, P{=}1)$, \\
        \textbf{horizon được công bố} (e.g.\ 100 vòng), \\
        đối thủ \textbf{khai thác được} (random, defection prob.\ 0.5--0.7).
    \end{block}
    \vspace{0.3cm}
    \begin{block}{Regime mà điều này đảo ngược}
        Horizon không xác định? Hợp tác toàn cầu hóa qua tất cả ngôn ngữ.\\
        Game cạnh tranh thuần (rock-paper-scissors)? LLM chơi \emph{sâu hơn} con người.
    \end{block}
    \vspace{0.3cm}
    \takeaway{Prosocial default là thật nhưng có điều kiện --- thuộc tính của game hợp tác \& mixed-motive, không phải trait đồng nhất của LLM.}
\end{frame}

\begin{frame}{Foreground mâu thuẫn --- bound claim một cách thành thật}
    \large
    \textbf{Bốn phát hiện đẩy ngược lại cách đọc mạnh về hiệu ứng ngôn ngữ:}
    \vspace{0.3cm}
    \normalsize
    \begin{enumerate}
        \setlength\itemsep{0.4em}
        \item \textbf{Persuasion null} (Biswas et al., CHI 2025): cross-lingual persuasive co-writing --- usage shifts, nhưng \emph{không có aggregate effect on persuasiveness}. Cái quan trọng là beliefs về AI authorship.
        \item \textbf{Cheap talk lấn át} (Madmoun \& Lahlou, 2025): một từ giao tiếp thêm nâng Stag Hunt cooperation từ \textbf{0\% lên 96.7\%} --- trong mọi ngôn ngữ.
        \item \textbf{Framing lấn át} (Buscemi et al., 2025): khi horizon không xác định, tỷ lệ hợp tác toàn cầu hóa qua các ngôn ngữ.
        \item \textbf{Vietnamese ``defection'' = lỗi số học} (Huynh et al., 2025): ``defection'' low-stakes đôi khi là payoff miscomprehension, không phải chiến lược.
    \end{enumerate}
    \vspace{0.3cm}
    \takeaway{Ngôn ngữ matters, nhưng không như núm xoay văn hóa cố định. Nó tương tác với framing, stake, giao tiếp, và task.}
\end{frame}

% =====================================================
% PHẦN 3 — TẠI SAO (CƠ CHẾ)
% =====================================================
\section{Tại sao --- cơ chế kiến trúc}

\begin{frame}{Hai cơ chế giải thích cross-lingual gap}
    \large
    \begin{columns}[T]
        \begin{column}{0.55\textwidth}
            \textbf{Semantic Hub} \scriptsize (Wu et al., 2024)\\[4pt]
            \normalsize Multilingual inputs project vào một hub trung gian English-dominant.\\[6pt]
            \footnotesize\itshape Bằng chứng causal: vector điều khiển kích hoạt derived từ tiếng Anh shift output Tây Ban Nha/Trung không kém gì native steering. Hub causally engaged, không vestigial.
            \vspace{0.4cm}\\
            \textbf{Translation Barrier} \scriptsize (Bafna et al., 2025)\\[4pt]
            \normalsize Nhiều failures là lỗi \emph{realisation lớp cuối} --- mô hình reasoning đúng bên trong, fail khi project ra ngôn ngữ đích.
        \end{column}
        \begin{column}{0.42\textwidth}
            \begin{block}{Khủng hoảng calibration}
                \centering
                \large \textbf{Expected Calibration Error}\\[6pt]
                \begin{tabular}{lc}
                    \toprule
                    \textbf{Ngôn ngữ} & \textbf{ECE} \\
                    \midrule
                    Anh        & 4.6\% \\
                    Phi-Anh    & \textcolor{rsclaret}{\textbf{23.1\%}} \\
                    \bottomrule
                \end{tabular}\\[6pt]
                \scriptsize Ahuja et al.\ 2022 (MEGA / XGLM, BLOOM, PaLM, 15-bin equal-width)
            \end{block}
            \vspace{0.2cm}
            \takeaway{Mô hình \emph{tự tin sai} trong ngữ cảnh phi-Anh.}
        \end{column}
    \end{columns}
\end{frame}

\begin{frame}{Chúng tôi không claim độc quyền --- ba alternatives}
    \large
    \begin{enumerate}
        \setlength\itemsep{0.5em}
        \item \textbf{Vi mạch đặc thù theo ngôn ngữ} \scriptsize (Pires, Schlinger, Garrette 2019 trên mBERT)\\
        \normalsize Một số failures bền vững theo cặp ngôn ngữ. Gợi ý các subnetwork tách biệt một phần.
        \item \textbf{Subspace interference} \scriptsize (lineage XLM-R, mT5)\\
        \normalsize Scale làm dịu nhưng không loại bỏ nhiễu cross-lingual. Low-resource langs trôi.
        \item \textbf{Latent-language priors} \scriptsize (Wendler et al., 2024)\\
        \normalsize Llama decoders đi qua intermediate \emph{gần với tiếng Anh hơn} ngôn ngữ đầu vào.
    \end{enumerate}
    \vspace{0.3cm}
    \begin{block}{Phép thử tự nhiên --- chạy được \emph{ngay}}
        \textbf{Mô hình ưu tiên ngôn ngữ}: CT-LLM (800B Trung) và TigerLLM (Bangla).\\[4pt]
        Nếu Semantic Hub là về \emph{dữ liệu huấn luyện} không phải kiến trúc, các mô hình này nên \emph{đảo ngược} sự bất đối xứng.
    \end{block}
\end{frame}

\begin{frame}{Đóng vòng lặp --- từ representation tới behavior}
    \large
    Liên kết từ ``hub asymmetry'' tới ``phân phối chiến lược IPD Tiếng Việt'' hiện tại là \textbf{inferential, không causal}.
    \vspace{0.3cm}
    \begin{block}{Tại sao điều này có thể fix bây giờ}
        Interpretability gần đây có đủ precedent cần thiết:
        \begin{itemize}
            \setlength\itemsep{0.2em}
            \item Spatial reps causally affect prediction \scriptsize (Chen et al., 2023)
            \item World-model state vars linearly decodable \emph{và} causally engaged \scriptsize (Zhang, 2026)
            \item Layer-resolved tools (LogitLens4LLMs, arithmetic-trajectory probes, syntax/semantics probes)
        \end{itemize}
    \end{block}
    \vspace{0.3cm}
    \takeaway{Thí nghiệm kế tiếp: derive steering vector từ \emph{Tiếng Anh} cooperative vs.\ defective IPD play. Apply vào cùng mô hình ở Tiếng Ả Rập. Đo strategy shift kết quả.\\Chuyển representation claim thành behavioural lever.}
\end{frame}

% =====================================================
% PHẦN 4 — TOM VÀ MODERATORS
% =====================================================
\section{Theory of Mind qua các ngôn ngữ}

\begin{frame}{Trẻ song ngữ giỏi ToM \emph{hơn}. LLM thì kém hơn.}
    \large
    \textbf{Baseline ở người:} trẻ song ngữ vượt qua bài false-belief \emph{sớm hơn} đơn ngữ.
    \scriptsize \itshape (Kovács 2009; Goetz 2003 --- đòi hỏi executive của code-switching scaffold ToM.)
    \vspace{0.3cm}
    \normalsize\textbf{Dự đoán ngầm cho máy:} mô hình đa ngôn ngữ nên là tác tử ToM \emph{mạnh hơn}, không yếu hơn.
    \vspace{0.4cm}
    \begin{block}{Cái ta thực sự thấy}
        \begin{itemize}
            \setlength\itemsep{0.3em}
            \item \textbf{XToM} (Chan et al., 2025): GPT-4o có khoảng cách ToM 11 điểm giữa Tiếng Nhật và Tiếng Trung trên cùng items.
            \item \textbf{LLM-Hanabi}: khi first-order ToM erodes, cả \emph{hợp tác} và \emph{lừa dối} cùng erode.
        \end{itemize}
    \end{block}
    \vspace{0.2cm}
    \footnotesize\itshape Caveat (Ullman 2023; Sap et al.\ 2022): điểm ToM tuyệt đối của LLM đang tranh cãi. Nhưng \emph{differences} cross-lingual trên parallel items vẫn well-defined bất kể construct-validity disputes.
\end{frame}

\begin{frame}{Ngôn ngữ là một moderator. Có ít nhất ba moderators khác.}
    \large
    \begin{columns}[T]
        \begin{column}{0.31\textwidth}
            \begin{block}{Personality steering}
                \footnotesize Ong et al.\ (2025).\\
                Steering Big-Five Agreeableness điều biến IPD cooperation \emph{và} tăng exploitation susceptibility.\\[4pt]
                \textit{$\Rightarrow$ Internal reps causally steer cooperation.}
            \end{block}
        \end{column}
        \begin{column}{0.31\textwidth}
            \begin{block}{Communication delay}
                \footnotesize Nishimoto et al.\ (2026).\\
                Quan hệ U-shape giữa delay và hợp tác. Agents \emph{tự phát} khai thác đối tác chậm hơn.\\[4pt]
                \textit{$\Rightarrow$ Hạ tầng đổi behavior, không cần đổi ngôn ngữ.}
            \end{block}
        \end{column}
        \begin{column}{0.31\textwidth}
            \begin{block}{Partner identity}
                \footnotesize Jiang et al.\ (2025).\\
                Hợp tác của con người chống agent phụ thuộc declared identity (người/rule-AI/LLM).\\[4pt]
                \textit{$\Rightarrow$ Ngôn ngữ là một signal định khung identity.}
            \end{block}
        \end{column}
    \end{columns}
    \vspace{0.4cm}
    \takeaway{Bất cứ phép thử linguistic relativity sạch nào phải hold ba cái này constant. Nếu không, ta nhầm moderator effects với hiệu ứng ngôn ngữ.}
\end{frame}

% =====================================================
% PHẦN 5 — KHUNG TIẾN HÓA
% =====================================================
\section{Hợp tác --- tiến hóa hay kế thừa?}

\begin{frame}{Hợp tác LLM không vừa khuôn Axelrod}
    \large
    \textbf{Truyền thống Axelrod:} hợp tác là một \emph{thành tựu}.\\
    Hợp tác ổn định đòi hỏi tương tác lặp lại, danh tiếng, hình phạt, network reciprocity.
    \vspace{0.3cm}
    \begin{block}{LLM đến trong trạng thái \emph{pre-cooperative}}
        Trước khi bất cứ game nào diễn ra:
        \begin{itemize}
            \setlength\itemsep{0.2em}
            \item \textbf{Pretraining corpora} = sản phẩm văn hóa tích lũy của hàng triệu authors người
            \item \textbf{RLHF} = một lớp ưu tiên người được chọn lọc
        \end{itemize}
        \vspace{0.2cm}
        Cùng nhau hoạt động như \textbf{kế thừa văn hóa} (Boyd \& Richerson 1985; Henrich 2015).
    \end{block}
    \vspace{0.3cm}
    \takeaway{Hợp tác LLM là \emph{kế thừa văn hóa}, không phải tiến hóa qua tương tác. Hệ quả là có thể kiểm tra.}
\end{frame}

\begin{frame}{Vấn đề WEIRD và transmission bottleneck đa ngôn ngữ}
    \large
    \begin{columns}[T]
        \begin{column}{0.50\textwidth}
            \textbf{Corpora WEIRD-skewed}\\
            \footnotesize Henrich, Heine, Norenzayan (2010); Muthukrishna et al.\ (2020):\\[4pt]
            \normalsize Quần thể WEIRD là \textbf{outlier thống kê} về individualism, tư duy phân tích, fair-procedure framings.
            \vspace{0.3cm}
            Nếu LLM kế thừa hợp tác từ corpora này, chúng encode chuẩn WEIRD trong khi \emph{xuất hiện như universal}.
            \vspace{0.3cm}
            \takeaway{Triển khai đa ngôn ngữ rủi ro \emph{áp đặt normative} --- không chỉ khoảng cách performance.}
        \end{column}
        \begin{column}{0.46\textwidth}
            \begin{alertblock}{Máy giờ đối mặt tiến hóa văn hóa}
                \footnotesize Quần thể LLM được liên kết qua:
                \begin{itemize}
                    \item Benchmark chia sẻ
                    \item Họ mô hình chia sẻ
                    \item Distillation \& self-training trên outputs của nhau
                \end{itemize}
                \vspace{0.2cm}
                Nếu cổ chai là Tiếng Anh (lingua franca của benchmark + RLHF annotators), \textbf{trait phi-Anh trôi độc lập}.
                \vspace{0.2cm}
                Tín hiệu thực nghiệm: khoảng cách English-Arabic trong TCG-Bench tại 32B params \emph{không} sửa được bằng scale.
            \end{alertblock}
        \end{column}
    \end{columns}
\end{frame}

% =====================================================
% PHẦN 6 — XÃ HỘI VÀ GOVERNANCE
% =====================================================
\section{Xã hội và governance}

\begin{frame}{Ba hệ quả xã hội --- đã đang triển khai}
    \large
    \begin{columns}[T]
        \begin{column}{0.32\textwidth}
            \begin{alertblock}{Chất lượng thể chế bất bình đẳng}
                \footnotesize Calibration \emph{kém hơn $\sim 5\times$} ngoài Tiếng Anh.\\[4pt]
                Public-facing agents ở Global South sẽ đưa ra lời khuyên \emph{tự tin sai} cho những người có ít lựa chọn nhất.\\[4pt]
                Đã triển khai: UlizaLlama (sức khỏe, Swahili), Aya Expanse (23 langs incl.\ VI), Bhashini (India).
            \end{alertblock}
        \end{column}
        \begin{column}{0.32\textwidth}
            \begin{alertblock}{Bề mặt tấn công bất đối xứng}
                \footnotesize Code-Switching Red-Teaming (Yoo et al., 2025): \textbf{+46.7\%} tỷ lệ tấn công thành công bằng cách trộn ngôn ngữ ở mức token.\\[4pt]
                Tấn công Agent-in-the-Middle khai thác kênh inter-agent.\\[4pt]
                Đối thủ giành đòn bẩy ở chính nơi safety data thưa nhất.
            \end{alertblock}
        \end{column}
        \begin{column}{0.32\textwidth}
            \begin{alertblock}{Đồng tiến hóa}
                \footnotesize Outputs LLM tái nhập diễn ngôn người, rồi vào training corpora tương lai.\\[4pt]
                Defaults máy điều biến theo ngôn ngữ định hình text mà người viết tiếp.\\[4pt]
                Vòng lặp kế thừa văn hóa thành \textbf{hai chiều}.
            \end{alertblock}
        \end{column}
    \end{columns}
\end{frame}

\begin{frame}{Ba đòn bẩy governance --- không cần đột phá}
    \large
    \begin{enumerate}
        \setlength\itemsep{0.5em}
        \item \textbf{Đa dạng hóa benchmark.}\\
        \normalsize Ghép benchmark an toàn (M-ALERT) với benchmark interactive multi-agent (clembench, MultiAgentBench, TCG-Bench). Báo cáo khoảng cách hành vi \emph{và} an toàn \emph{cùng nhau}.

        \item \textbf{Chuẩn coverage RLHF đa ngôn ngữ.}\\
        \normalsize Mọi model card báo cáo annotator composition theo ngôn ngữ và phương ngữ. Đề xuất floor: $\geq 5\%$ preference labels per language cho bất cứ model triển khai qua $> 5$ ngôn ngữ, với dialectal stratification được tài liệu hóa.

        \item \textbf{Chuẩn audit cho triển khai multi-agent.}\\
        \normalsize Probe code-switching và mô hình mối đe dọa Agent-in-the-Middle thành test bed \emph{mặc định} --- như probe prompt-injection đã là cho single-agent.
    \end{enumerate}
    \vspace{0.3cm}
    \takeaway{Cái thiếu là \emph{hạ tầng thể chế}, không phải methodology.}
\end{frame}

% =====================================================
% PHẦN 7 — RESEARCH AGENDA
% =====================================================
\section{Bốn dự đoán có thể bác bỏ}

\begin{frame}{Bốn dự đoán, mỗi cái falsifiable với phương pháp hiện tại}
    \begin{columns}[T]
        \begin{column}{0.48\textwidth}
            \begin{block}{P1 --- Corpus $\leftrightarrow$ hợp tác}
                \footnotesize Mật độ cooperative-discourse trong corpora tiền huấn luyện $\rho$ tỷ lệ hợp tác IPD qua các ngôn ngữ.\\
                Workarounds cho closed corpora: open-data models (Pile, CT-LLM, TigerLLM), proxy corpora, controlled small-model runs.\\[4pt]
                \textbf{Falsifier:} Spearman $\rho$ không phân biệt được khỏi 0 qua $\geq 8$ ngôn ngữ.
            \end{block}
            \vspace{0.2cm}
            \begin{block}{P2 --- Hub probe + behavioural readout}
                \footnotesize Probe layers $\lfloor L/3 \rfloor, \lfloor 2L/3 \rfloor$ phân loại hub vs.\ realisation failures.\\
                \textbf{Sau đó:} áp dụng English-derived steering vector vào IPD prompt phi-Anh; đo strategy shift.\\[4pt]
                \textbf{Falsifier:} không có behavioural shift sau steering.
            \end{block}
        \end{column}
        \begin{column}{0.48\textwidth}
            \begin{block}{P3 --- Drift dưới distillation}
                \footnotesize Chuỗi distillation kiểm soát: traits truyền dẫn high-fidelity trong Tiếng Anh, suy giảm trong low-resource langs.\\[4pt]
                \textbf{Falsifier:} không có degradation $\Rightarrow$ inheritance-via-distillation sai.
            \end{block}
            \vspace{0.2cm}
            \begin{block}{P4 --- Ceiling độ tin cậy đánh giá}
                \footnotesize Multilingual LLM-as-Judge với $\geq 3$ judges dị thể + self-translation + checklist rubrics nâng Fleiss' $\kappa > 0.5$.\\[4pt]
                \textbf{Falsifier:} saturates dưới $0.5$ $\Rightarrow$ cần bilingual human adjudication, không phải aggregation tốt hơn.
            \end{block}
        \end{column}
    \end{columns}
\end{frame}

% =====================================================
% PHẦN 8 — KẾT LUẬN
% =====================================================
\section{Take-aways}

\begin{frame}{Năm điều cần nhớ}
    \large
    \begin{enumerate}
        \setlength\itemsep{0.5em}
        \item \textbf{Ngôn ngữ điều biến machine behavior.} Là biến chủ động, không phải phương tiện thụ động.
        \item \textbf{Cơ chế là kiến trúc.} Hub, Barrier, microcircuits, latent-language --- không phải cultural essentialism.
        \item \textbf{Hợp tác là kế thừa, không phải tiến hóa.} Pretraining + RLHF là cơ chế kế thừa.
        \item \textbf{Mâu thuẫn là informative.} Persuasion null, cheap-talk dominance, lỗi số học Tiếng Việt --- chúng bound claim.
        \item \textbf{Tổng hợp $\rightarrow$ research agenda.} Bốn dự đoán falsifiable; đóng vòng lặp representation--behavior là cấp bách nhất.
    \end{enumerate}
\end{frame}

\begin{frame}{Vị trí trong chương trình machine-behaviour}
    \large
    \textbf{Liên tục với:}\\
    \footnotesize Rahwan (machine behaviour) $\cdot$ Han (evolutionary game theory cho multi-agent emergence) $\cdot$ Perc (statistical physics of cooperation) $\cdot$ Boyd, Richerson, Henrich (cultural evolution) $\cdot$ Axelrod, Nowak (cooperation without design)
    \vspace{0.5cm}\\
    \normalsize\textbf{Đóng góp riêng biệt:}
    \begin{itemize}
        \setlength\itemsep{0.3em}
        \item Treats \emph{ngôn ngữ} như biến điều biến trong machine ethology
        \item Tái định khung ``nicer than human'' như kế thừa văn hóa, với hệ quả WEIRD
        \item Bridges mechanistic interpretability với population-level cooperation theory
        \item Cung cấp agenda falsifiable ngay
    \end{itemize}
    \vspace{0.5cm}
    \takeaway{Quần thể máy giờ cung cấp một substrate thực nghiệm mới --- consequential một cách khó chịu --- cho khoa học về hợp tác.}
\end{frame}

\begin{frame}
    \begin{center}
        \Huge \textcolor{rsnavy}{\textbf{Cảm ơn}} \\[0.6cm]
        \large Câu hỏi \& thảo luận \\[0.8cm]
        \normalsize
        \textbf{Đào Sỹ Duy Minh}\\
        Khoa Công nghệ Thông tin, Trường Đại học Khoa học Tự nhiên, ĐHQG-HCM\\
        \texttt{dsdminh22@apcs.fitus.edu.vn}\\[0.8cm]
        \scriptsize\itshape Babel's Machines --- Linguistic Relativity, Cooperation,\\and the Emerging Behaviour of Multilingual Artificial Agents.\\
        Cho Interface Focus theme issue về machine behaviour in the age of LLMs.
    \end{center}
\end{frame}

\end{document}
